% Reusable 16.5 cm wide figure; requires amsmath, amssymb, TikZ,
% and the TikZ libraries arrows.meta and calc.
\definecolor{tlBlue}{HTML}{2463AD}
\definecolor{tlGold}{HTML}{A46900}
\definecolor{tlPink}{HTML}{9D3D86}
\definecolor{tlGreen}{HTML}{23734F}
\definecolor{tlInk}{HTML}{202B35}
\definecolor{tlMute}{HTML}{89929B}
\definecolor{tlFaint}{HTML}{E1E6EA}
\newcommand{\ThreeLineNonuniquenessFigureVThree}{%
\begin{tikzpicture}[x=1cm,y=1cm,text=tlInk,
  font=\fontsize{11}{15}\selectfont,
  tl heading/.style={anchor=west,font=\sffamily\bfseries\fontsize{12.5}{15}\selectfont},
  tl label/.style={fill=white,inner sep=2pt},
  tl step/.style={-{Latex[length=3mm,width=2mm]},line width=1.65pt},
  tl zero/.style={circle,draw=tlMute,fill=white,minimum size=3pt,inner sep=0pt}]
  \path[use as bounding box] (0,0) rectangle (16.5,22.6);
  \node[anchor=west,font=\sffamily\bfseries\fontsize{17}{20}\selectfont]
    at (0,22.20) {Do three different lines give unique summands?};
  \node[anchor=west,font=\fontsize{12}{15}\selectfont] at (0,21.40)
    {$\textcolor{tlBlue}{U_1=\{(x,0)\}},\qquad
      \textcolor{tlGold}{U_2=\{(0,y)\}},\qquad
      \textcolor{tlPink}{U_3=\{(t,t)\}}$\quad over $\mathbb R$.};
  \node[anchor=west,font=\fontsize{11.5}{15}\selectfont] at (0,20.65)
    {Each pair meets only at $0$. Keep the target $\textcolor{tlGreen}{v=(1,1)}$ fixed.};
  \draw[tlFaint,line width=.5pt] (0,13.75)--(16.5,13.75);
  \draw[tlFaint,line width=.5pt] (0,7.15)--(16.5,7.15);
  \foreach \py/\step/\tt/\heading in {
    13.9/1/0/{Start with $u_3=0$.},
    7.3/2/.5/{Use $u_3=\left(\frac12,\frac12\right)$.},
    .7/3/1/{Use $u_3=(1,1)$.}}{
    \begin{scope}[shift={(0,\py)}]
      \node[tl heading] at (0,5.75) {\step\quad\heading};
      \pgfmathsetmacro{\remaining}{1-\tt}
      \begin{scope}[shift={(1.30,1.60)},x=3cm,y=3cm]
        % Exactly the same camera, subspaces, target, and scale in each row.
        \draw[tlBlue!35,line width=1.2pt] (-.2,0)--(1.42,0);
        \draw[tlGold!35,line width=1.2pt] (0,-.2)--(0,1.22);
        \draw[tlPink!35,line width=1.2pt] (-.2,-.2)--(1.28,1.28);
        \node[tl label,text=tlBlue,anchor=west] at (1.42,0) {$U_1$};
        \node[tl label,text=tlGold,anchor=east] at (-.07,1.15) {$U_2$};
        \node[tl label,text=tlPink,anchor=west] at (1.23,1.25) {$U_3$};
        \node[below left,text=tlMute,font=\small] at (0,0) {$0$};
        \fill[tlMute] (0,0) circle[radius=1.5pt];
        \ifnum\step<3
          \draw[tl step,tlBlue] (0,0)--(\remaining,0);
          \node[tl label,text=tlBlue,anchor=north] at ({.5*\remaining},-.045) {$u_1$};
          \draw[tl step,tlGold] (\remaining,0)--(\remaining,\remaining);
          \node[tl label,text=tlGold,anchor=east]
            at ({\remaining-.055},{.5*\remaining}) {$u_2$};
        \fi
        \ifnum\step>1
          \draw[tl step,tlPink] (\remaining,\remaining)--(1,1);
          \node[tl label,text=tlPink,anchor=south east]
            at ({.5*(1+\remaining)-.025},{.5*(1+\remaining)+.035}) {$u_3$};
        \fi
        \ifnum\step=2
          \fill[tlGold] (.5,.5) circle[radius=2.1pt];
          \node[tl label,anchor=south east,font=\fontsize{10.5}{13}\selectfont]
            at (.46,.55) {$\left(\frac12,\frac12\right)$};
        \fi
        \fill[tlGreen] (1,1) circle[radius=3pt];
        \node[tl label,text=tlGreen,anchor=west] at (1.055,1.01) {$v=(1,1)$};
      \end{scope}
      \ifnum\step=1
        \node[anchor=north west,align=left,font=\fontsize{13}{18}\selectfont] at (8.30,4.70)
          {$\textcolor{tlBlue}{u_1=(1,0)}$\\[6pt]
           $\textcolor{tlGold}{u_2=(0,1)}$\\[6pt]
           $\textcolor{tlPink}{u_3=(0,0)}$};
        \node[anchor=west] at (8.30,1.35) {The two axis vectors already add to $v$.};
      \fi
      \ifnum\step=2
        \node[anchor=north west,align=left,font=\fontsize{13}{18}\selectfont] at (8.30,4.95)
          {$\textcolor{tlBlue}{u_1=\left(\frac12,0\right)}$\\[6pt]
           $\textcolor{tlGold}{u_2=\left(0,\frac12\right)}$\\[6pt]
           $\textcolor{tlPink}{u_3=\left(\frac12,\frac12\right)}$};
        \node[anchor=west,align=left] at (8.30,.95)
          {The axis vectors add to $\left(\frac12,\frac12\right)$.\\[4pt]
           Adding the diagonal vector still gives $v$.};
      \fi
      \ifnum\step=3
        \node[anchor=north west,align=left,font=\fontsize{13}{18}\selectfont] at (8.30,4.70)
          {$\textcolor{tlBlue}{u_1=(0,0)}$\\[6pt]
           $\textcolor{tlGold}{u_2=(0,0)}$\\[6pt]
           $\textcolor{tlPink}{u_3=(1,1)}$};
        \node[anchor=west,align=left] at (8.30,1.05)
          {The diagonal vector alone equals $v$.\\[6pt]
           The summands changed; $v$ did not.};
      \fi
    \end{scope}
  }
  \node[anchor=west,font=\sffamily\bfseries\fontsize{12}{15}\selectfont] at (0,.25)
    {The same vector has different summands. This sum is not direct.};
\end{tikzpicture}%
}
